\lyseis{7.4}{Αφαίρεση ρητών}
\sol{1} α) $2$ \quad β) $-2$ \quad γ) $-3$ \quad δ) $3$ \quad ε) $-6$ \quad στ) $6$.
\sol{2} α) $7+(-9)=-2$ \quad β) $-4+6=2$ \quad γ) $3{,}5+1{,}5=5$ \quad δ) $-\frac{1}{2}-\frac{1}{3}=-\frac{5}{6}$ \quad ε) $-2{,}4-0{,}6=-3$ \quad στ) $8+8=16$.
\sol{3} α) $x=3-7=-4$ \quad β) $x=-2+5=3$ \quad γ) $x=4-9=-5$ \quad δ) $x=-6-(-6)=0$ \quad ε) $x=-3-2=-5$.
\sol{4} α) $10-3+8-2=13$ \quad β) $5-3+7-2+9=16$ \quad γ) $4-5-7+10=2$ \quad δ) $20+5+3=28$.
\sol{5} α) $-12$ \quad β) $-2{,}5+7-4{,}5-1=-1$ \quad γ) $\frac{1}{2}-\frac{3}{4}+\frac{1}{4}-1=-1$.
\sol{6} α) Λάθος. β) Σωστό. γ) Σωστό. δ) Σωστό. ε) Λάθος: $4+2=6$. στ) Σωστό ($3-5=-2$).
\sol{7} α) $5-(-3)=8$ \quad β) $-2-(-11)=-2+11=9$.
\sol{8} $2\,917-(-430)=2\,917+430=3\,347$ μ.
\sol{9} $250-(-206)=456$ χρόνια (περίπου, χωρίς έτος $0$: $455$).
\sol{10} $8-(-15)=23$ € (ή $-15-8=-23$ αν αφαιρέσουμε ανάποδα).
\sol{11} $(-9)-(-22)=13$ μ.
\sol{12} $(-2)-3=-5$: κατεβαίνει $5$ ορόφους.
\sol{13} $-7-4=-11$: πέφτει $11$ βαθμούς.
\sol{14} Το σωστό είναι $5-(-3)=5+3=8$. Στον άξονα από το $-3$ ως το $5$ είναι $8$ μονάδες. Ο Στέλιος αφαίρεσε το $3$ αντί για το $-3$.
\sol{15} Όχι: $3-(-2)=5>3$.
\sol{16} Θετικό αν $\alpha>\beta$, αρνητικό αν $\alpha<\beta$, μηδέν αν ίσα.
\sol{17} $10-(3+4)=3$ και $10-3-4=3$.
\sol{18} $x-(-5)=2$, άρα $x=2+(-5)=-3$.
\sol{19} Ελεύθερη απάντηση.
\sol{20} Ελεύθερη απάντηση.
